% ======================================================================
% Kapitel 6 – Tabellen (AP2 laut review/plan-kap6.md, Abschnitt 9)
% Entwurf, noch nicht in Kapitel/06_Konzept.tex eingebunden (AP8 fügt zusammen).
% Code-Stand aller Zellen: G1@25d5c1a (11.11.2026). Pfade in den Kommentaren
% sind relativ zu g1pilot/ (Doku: docs/, Code: g1pilot/..., launch/...).
% Stil: Gitter nach CLAUDE.md (\tabstil, tabularx, L/C-Spalten, \hline je Zeile,
% \tabkopfzeile + \tabkopf{}, Caption oben). Eigenleistung: kein \cite.
% Tabellen enthalten keine Dateinamen, Knotennamen oder Tastennamen
% (Abgrenzung 6/7, Plan Abschnitt 2); Bedienschritte stehen im Anhang~\ref{anh:dokumentation}.
% Texteinleitung zu jeder Tabelle (\ref vor dem Float) schreibt AP3 bis AP7.
% ======================================================================


% ----------------------------------------------------------------------
% Tabelle 1: Abgrenzung der Eigenleistung (Abschnitt 6.1.3)
% ----------------------------------------------------------------------
\begin{table}[htbp]
    \tabstil
    \centering
    \caption[Abgrenzung der Eigenleistung]{Abgrenzung der Eigenleistung nach Funktionsbereichen (übernommen, erweitert, neu). Prozentangaben sind Näherungswerte der Herkunftsangabe im Paket je Programmdatei.}
    \label{tab:eigenleistung}
    \begin{tabularx}{\textwidth}{|>{\raggedright\arraybackslash}p{3.7cm}|>{\raggedright\arraybackslash}p{2.5cm}|L|}\hline
        \tabkopfzeile \tabkopf{Funktionsbereich} & \tabkopf{Stufe} & \tabkopf{Bemerkung}\\\hline
        Roboterzustand & übernommen & Anteil des Originals etwa 80\,\%.\\\hline
        Navigation (Simulation) & übernommen, eingebunden & Originalanteile 35 bis 80\,\% je Datei; läuft nur in der Simulation.\\\hline
        Armregelung und inverse Kinematik & erweitert & Original: Grundgerüst (etwa 20\,\%), Löser (etwa 30\,\%), Zielmarker (etwa 50\,\%). Ergänzt: Überblendung, Begrenzer.\\\hline
        Anbindung des Onboard-Reglers & erweitert & Original: Grundgerüst (etwa 20\,\%). Ergänzt: Bodenkontaktprüfung, Halt-Wege.\\\hline
        Kollisionsmodell und Bahnplaner & neu & Kugelmodell, Laufzeitprüfung, eigener Bahnplaner.\\\hline
        Simulationskopplung & neu & Basis: Simulationspaket des Herstellers; Kopplung, Zeitführung und Szenen eigen. \pruefen{Herkunft und Lizenz des Szenen-Editors (mujoco-scene-editor) in der Herkunftsangabe des Pakets ergänzen}\\\hline
        Ersatzregler Fortbewegung & neu, Policies übernommen & Standregler und Einbindung neu; Policies von Unitree (eine) und NVIDIA (zwei).\\\hline
        Sicherheitsfunktionen & neu & Tabelle~\ref{tab:haltstufen} und Tabelle~\ref{tab:schutzmechanismen}.\\\hline
        Hände (Inspire, eigene Profile) & neu & Inspire-Hände in drei Stufen, Hand-Profile für die Simulation; Dex3-Hand des Originals entfernt. \pruefen{Herkunft der Referenzskripte für die Modbus-Register klären}\\\hline
        Bedienoberfläche, Gamepad, Smartphone & neu & Aus dem Original: Basisklasse, Joystick-Eingang, Umschalter (je 25 bis 80\,\%).\\\hline
        Posen und Abläufe & neu & Posenspeicher, relative Bewegungen, Abläufe, Editor.\\\hline
        Arm-API & neu & HTTP-Schnittstelle mit JSON, ohne ROS nutzbar.\\\hline
        Bereitstellung und Startmenü & neu & Container-Images und Profile, Startmenü; großes Docker-Image des Originals entfernt.\\\hline
        Tests und Dokumentation & neu & 597 Testfunktionen (rund 13.400 Zeilen); Paket rund 26.400 Zeilen Python, davon etwa ein Zehntel Original. Dokumentation: Anhang~\ref{anh:dokumentation}.\\\hline
    \end{tabularx}
\end{table}
% Fundstellen je Zeile (Stand G1@25d5c1a, Pfade relativ zu g1pilot/):
% Zählung Zeile 14 neu am 20.11.2026: find g1pilot -name "*.py" | wc -l = 26.398 Zeilen; test/*.py = 13.367 Zeilen;
%   grep "def test_" in test/ = 597 Funktionen (deckt sich mit der Zählung im Plan vom 16.11.).
% 1 Roboterzustand:      NOTICE.md Tabelle: g1pilot/state/robot_state.py "~80 %"; docs/02-T-Architektur.md "robot_state" (Gelenke, IMU, TF).
% 2 Navigation:          NOTICE.md: navigation/dijkstra_planner.py, nav2point.py, create_map.py "35–80 %"; docs/31-A-Navigation.md "läuft nur in der Simulation",
%                        Pose aus MuJoCo-Ground-Truth (sim_localization.py, nicht in NOTICE-Tabelle = eigen), scene_bridge.py (Hindernisse, eigen).
% 3 Armregelung/IK:      NOTICE.md: arm_controller.py "Grundgerüst (~20 %)", utils/ik_solver.py "~30 %", manipulation/interactive_marker.py "~50 %";
%                        eigen: manipulation/ik_servo_process.py (nicht in NOTICE), Gewichtsrampe arm_controller.py:310-311, REAL_MAX_* :1119-1124.
% 4 Onboard-Anbindung:   NOTICE.md: navigation/loco_client.py "Grundgerüst (~20 %)"; eigen: navigation/ground_contact.py, loco_client.py:_halt_legs (:438),
%                        Zustandsabgleich state_sync_hz (:163, :457).
% 5 Kollision/Planer:    manipulation/arm_geometry.py, arm_planner.py, arm_planner_process.py (nicht in NOTICE = eigen); docs/20-T-Arm-Manipulation.md
%                        "arm_planner.py: Eigenes System (kein MoveIt/OMPL)".
% 6 Simulationskopplung: NOTICE.md "unitree_mujoco/ ... Unitree Robotics BSD-3-Clause"; docs/02-T-Architektur.md "MuJoCo-Simulation (Unitree-Original, angepasst)";
%                        plan-kap6.md N3 (Nils 16.11.): Original nur Realbetrieb, alles Simulationsbezogene Nils' Werk.
%                        Szenen-Editor: unitree_mujoco/scene_editor/README.md "Grundlage ist der mujoco-scene-editor"; in NOTICE.md nicht aufgeführt -> \pruefen.
% 7 Ersatzregler:        NOTICE.md "policies/g1_wholebody/: Unitree unitree_rl_mjlab, Apache-2.0", "policies/agile_velocity_*: NVIDIA WBC-AGILE";
%                        docs/30-T-Stehen-und-Gehen.md "Zwei kombinierte Regler"; loco_sim.py, stand_balancer.py, com_model.py nicht in NOTICE = eigen.
% 8 Sicherheit:          siehe Tabellen tab:haltstufen und tab:schutzmechanismen.
% 9 Hände:               docs/22-T-Inspire-Haende.md "Die drei Backend-Stufen"; docs/25-T-Eigene-Haende.md; NOTICE.md "Entfernt: Dex3-Hand";
%                        ../reference/ftp_hand_controller/README.md "vom Branch reference/ftp_hand_controll uebernommen, Commit 78337ea" (Urheber unklar -> \pruefen).
% 10 Bedienwege:         NOTICE.md: teleoperation/joy_mux.py, joystick.py, gui_base.py "25–80 %"; "Entfernt: Streamdeck-GUI";
%                        eigen: demo_gui.py, ps4_*.py, phone/ (NOTICE "Alle übrigen Dateien ... Demo-GUI, PS4/Xbox-Teleop, Handy-Zugang").
% 11 Posen/Abläufe:      manipulation/pose_store.py, teleoperation/sequence_runner.py, sequence_editor.py (nicht in NOTICE = eigen); docs/23-T-Posen-und-Ablaeufe.md.
% 12 Arm-API:            manipulation/arm_api.py, arm_command.py (NOTICE: "Arm-API ... eigene Arbeit"); docs/21-T-Arm-API.md.
% 13 Bereitstellung:     docker-compose.yml (Images g1pilot-mujoco, g1pilot-sim, g1pilot-real), docker/Dockerfile.{mujoco,sim,real}; start.sh, g1_gui.py;
%                        NOTICE.md "Entfernt: ... das große Docker-Image".
% 14 Tests/Doku:         siehe Zählung oben; NOTICE.md "Stand 2026-10-30 stammt nur noch etwa ein Zehntel der Zeilen aus dem Original"; docs/ (45 Dateien).


% ----------------------------------------------------------------------
% Tabelle 2: Unterschiede Simulation / Realbetrieb (Abschnitt 6.2.2)
% ----------------------------------------------------------------------
\begin{table}[htbp]
    \tabstil
    \centering
    \caption[Unterschiede zwischen Simulation und Realbetrieb]{Unterschiede zwischen Simulation und Realbetrieb. Die Befehle der Bedienwege und der Funktionsknoten sind in beiden Fällen dieselben.}
    \label{tab:sim_real}
    \begin{tabularx}{\textwidth}{|>{\raggedright\arraybackslash}p{2.3cm}|L|L|L|}\hline
        \tabkopfzeile \tabkopf{Aspekt} & \tabkopf{Simulation} & \tabkopf{Realbetrieb} & \tabkopf{gleich?}\\\hline
        Verbindung & DDS-Domain 1 über Loopback & Domain 0 über die Netzwerkkarte & Nachrichten ja, Domain nein (ein Schalter)\\\hline
        Fortbewegung & Ersatzregler aus Standregler und Lauf-Policy & Onboard-Regler des Roboters & Befehle ja, Regler nein (Abschnitt~\ref{sec:fortbewegungspfade})\\\hline
        Arme & Umschalten ohne Rampe, keine Rückfahrt & Rampe und geplante Rückfahrt vor der Rückgabe & Regelung ja, Übergabe nein\\\hline
        Hände & simulierte Finger; Inspire-Hand, Gummihand oder eigenes Profil & Inspire-Hände über Modbus TCP & Schnittstelle ja, Anbindung nein\\\hline
        Grenzwerte & Voreinstellungen, nicht durchgesetzt & durchgesetzt (Tabelle~\ref{tab:schutzmechanismen}) & nein\\\hline
        Start & automatischer Start in den Modus GREIFEN & kein selbsttätiger Start, Sicherheitsbestätigung & nein\\\hline
        Navigation & verfügbar, Karte und Position aus dem Simulator & nicht verfügbar & nein\\\hline
        Not-Halt (Beine) & Beine gedämpft, der Roboter fällt & Beine bleiben stehen (Tabelle~\ref{tab:haltstufen}) & nein\\\hline
    \end{tabularx}
\end{table}
% Fundstellen je Zeile (Doku / Code):
% Verbindung:   Doku docs/02-T-Architektur.md "Sim ↔ Real: Ein Schalter" (Domain 1/lo gegen 0/NIC); docs/04-T-Umgebungsvariablen.md G1_SIM_MODE.
%               Code g1pilot/utils/common.py:21-23 (SIM_DDS_DOMAIN=1, REAL_DDS_DOMAIN=0, SIM_DDS_INTERFACE="lo"), :36-40 resolve_dds, init_dds als einzige Stelle.
% Fortbewegung: Doku docs/30-T-Stehen-und-Gehen.md "Warum zwei völlig verschiedene Implementierungen?", "Zwei kombinierte Regler", Tabelle loco_client;
%               Code navigation/loco_sim.py (Zustände :91-98, Topics :251-254 gleich wie loco_client), navigation/loco_client.py.
% Arme:         Doku docs/20-T-Arm-Manipulation.md "rt/arm_sdk-Gewichtsrampe" ("Die Sim setzt beide Rampenzeiten auf 0.0 ... real 2.0 s"), "Übergabe ... (nur Echtbetrieb relevant)".
%               Code launch/bringup_sim.launch.py:109-110 (arm_weight_ramp_*_s = 0.0); manipulation_launcher.launch.py:74-75 (Default 2.0, real);
%               arm_controller.py:1349-1354 (_release_possible: Rückfahrt nur bei ramp > 0, Kommentar :318 "Sim (Rampe 0) unveraendert"), :1356 _start_release.
% Hände:        Doku docs/22-T-Inspire-Haende.md "Die drei Backend-Stufen" (Stufe 2 mujoco = Sim-Default, Stufe 3 modbus = Real-Default);
%               docs/25-T-Eigene-Haende.md "Eigene Profile wirken nur bei G1_SIM_MODE=true"; docs/04-T-Umgebungsvariablen.md G1_HANDS ("Real wirkungslos").
%               Profile inspire, rubber, eigenes: docs/04-T G1_HANDS.
% Grenzwerte:   Doku docs/50-A-Echter-Roboter.md E "Harte Grenzen nur auf real"; Code arm_controller.py:1029-1030 (if not is_sim_mode(): _apply_real_safety_limits),
%               :1126-1147; Sim höher: bringup_sim.launch.py:112 (ee_velocity_limit 0.5), manipulation_launcher.launch.py:53-56 (kp_low 150, kd_low 12 gegen real 60/1,5, bringup_real.launch.py:115-118).
% Start:        Doku docs/50-A-Echter-Roboter.md E "Kein Auto-Start", docs/40-T-Simulation.md "Auto-Start (GREIFEN)", docs/10-A-Demo-GUI.md "Auf dem echten Roboter sind die Kacheln gesperrt, bis Start gedrückt wurde";
%               Code teleoperation/sim_autostart.py (nur is_sim_mode), start.sh:197-198 und g1_gui.py:1143 (G1_REAL_CONFIRM).
% Navigation:   Doku docs/31-A-Navigation.md "Die Navigation läuft nur in der Simulation", docs/04-T G1_ENABLE_NAV (Default 0, "nur Sim");
%               Code navigation/sim_localization.py (Pose aus MuJoCo), scene_bridge.py.
% Not-Halt:     Doku docs/51-A-Notaus.md "In der Simulation" ("Der NOT-HALT setzt auch die Beine auf DAMP ... Voll-Notaus wirkt dort wie der NOT-HALT"), Tabelle "Stufen";
%               Code navigation/loco_sim.py:57 und :251 (_on_emergency -> DAMP), navigation/loco_client.py:438-454 (_halt_legs: StopMove, Damp nur mit estop_damp_legs), :157 (Default False).


% ----------------------------------------------------------------------
% Tabelle 3: Halt-Stufen (Abschnitt 6.3.1)
% Beschriftung der Auslöser nach Bedienweg, ohne Tastennamen (siehe Anhang).
% ----------------------------------------------------------------------
\begin{table}[htbp]
    \tabstil
    \setlength{\tabcolsep}{3pt}
    \centering
    \caption[Halt-Stufen]{Halt-Stufen mit Auslöser und Wirkung auf die Gelenkgruppen im Realbetrieb; \enquote{--} bedeutet unverändert. In der Simulation dämpft der Not-Halt zusätzlich die Beine (Tabelle~\ref{tab:sim_real}).}
    \label{tab:haltstufen}
    \begin{tabularx}{\textwidth}{|>{\raggedright\arraybackslash}p{1.9cm}|>{\hsize=1.2\hsize}L|>{\hsize=0.9\hsize}L|>{\hsize=0.9\hsize}L|>{\hsize=0.8\hsize}L|>{\hsize=1.2\hsize}L|>{\hsize=1.0\hsize}L|}\hline
        \tabkopfzeile \tabkopf{Stufe} & \tabkopf{Auslöser} & \tabkopf{Arme} & \tabkopf{Hände} & \tabkopf{Taille} & \tabkopf{Beine} & \tabkopf{Quittierung}\\\hline
        Stopp & Gamepad, Bedienoberfläche, Smartphone & Bewegung endet & -- & -- & -- & nicht nötig\\\hline
        Gehen stoppen & Totmann loslassen, Gehbefehl bleibt aus & -- & -- & -- & Roboter steht und balanciert & nicht nötig\\\hline
        Not-Halt & Bedienoberfläche, Gamepad, Smartphone & schlaff, gedämpft & halten Stellung & gehalten & Stopp, Onboard-Regler balanciert & Bedienoberfläche oder Smartphone\\\hline
        Voll-Notaus & Gamepad, 2\,s halten & schlaff, gedämpft & halten Stellung & passiv & gedämpft, Roboter sackt zusammen & nur Sperre; Beine bleiben gedämpft\\\hline
        Fern\-bedienung des Herstellers & Dämpfen an der Fernbedienung & \multicolumn{4}{p{7.7cm}|}{Dämpfung durch die Firmware, unabhängig von PC und Umgebung. \pruefen{Wirkung auf die Inspire-Hände klären}} & an der Fernbedienung\\\hline
    \end{tabularx}
\end{table}
% Fundstellen je Zeile (Doku / Code):
% Stopp:        Doku docs/51-A-Notaus.md Stufen-Tabelle ("Stopp ... Hand steht sofort, geplante Bewegung bricht ab. Roboter bleibt, wie er ist"),
%               docs/12-T-Handy.md Aktionstabelle (`cancel` = Bewegung stoppen), docs/10-A-Demo-GUI.md ("Bewegung stoppen").
%               Code teleoperation/ps4_arm_teleop.py (Kreuz = Stopp, Kopfkommentar :45-58), manipulation/arm_controller.py (pose_store/cancel).
%               "Hände unverändert" ist Schluss aus "Roboter bleibt, wie er ist" (Doku nennt Finger nicht) -> bei Bedarf mit Nils prüfen.
% Gehen stoppen: Doku docs/51-A-Notaus.md ("L2 loslassen (Totmann), GUI-Gehknopf loslassen ... G1 bleibt stehen, balanciert weiter"),
%               docs/10-A-Demo-GUI.md ("loslassen = sofort 0"); Code teleoperation/ps4_walk.py:47-52 (Deadman), navigation/loco_client.py:716-760 (_cmd_vel_tick, StopMove genau einmal), :148 cmd_vel_timeout.
% Not-Halt:     Doku docs/51-A-Notaus.md Stufen und "Quittieren" (Tabelle Sim/Echt: Demo-GUI Startleiste / Handy Start; Controller "geht nicht, absichtlich"),
%               docs/51-T-Notaus.md "Wer was tut", docs/50-A-Echter-Roboter.md "Was EMERGENCY STOP konkret bewirkt".
%               Code manipulation/arm_controller.py:1840-1850 (_publish_arm_slack, kp 0, kd = estop_arm_kd :221), Taille estop_waist_hold :227 (Default true),
%               launch/bringup_real.launch.py:128; navigation/loco_client.py:438-454 (_halt_legs, estop_damp_legs :157 False), manipulation/inspire_ftp/bridge.py:251-256 (set_enabled(False)).
%               Smartphone-Quittieren: teleoperation/phone/gateway.py:64 ("ack": Bediener, Steuerhoheit); Not-Halt auch in Rolle Ansicht: gateway.py:60 ("estop": Ansicht), link.py:68 estop_for_view.
% Voll-Notaus:  Doku docs/51-A-Notaus.md ("R1 + L1 + PS 2 s ... ganzer G1 in Damp ... sackt zusammen"; "Nach dem Voll-Notaus quittiert Start nur; der G1 bleibt in Damp"),
%               docs/51-T-Notaus.md (arm_controller: "wie NOT-HALT, zusätzlich Taille passiv"; loco_client: "Damp zuerst, dann StopMove").
%               Code teleoperation/ps4_arm_teleop.py:736 (full_estop_hold_s 2.0), arm_controller.py:1916 (_on_full_emergency_stop), loco_client.py:555-581 (Damp immer, Quittieren bleibt Damp).
%               Hände: der Voll-Notaus beginnt mit dem Not-Halt des PS-Drucks (51-A: "der NOT-HALT vom PS-Druck bleibt aber bestehen").
% Fernbedienung: Doku docs/51-A-Notaus.md ("Unitree-Fernbedienung: Damp an der Fernbedienung; wirkt an der Firmware, unabhängig von PC und Stack").
%               Wirkung auf die Inspire-Hände ist nirgends dokumentiert (Hände hängen per Modbus an eigenem Netz, docs/02-T-Architektur.md) -> \pruefen.


% ----------------------------------------------------------------------
% Tabelle 4: Schutzmechanismen mit Auslegungswerten (Abschnitt 6.3.2)
% Mehrseitig, daher xltabular. Grenzwerte stehen nur hier (Abgrenzung 6/7).
% ----------------------------------------------------------------------
\begingroup
\tabstil
\setlength{\tabcolsep}{3pt}
\begin{xltabular}{\textwidth}{|>{\hsize=1.1\hsize\raggedright\arraybackslash}X|>{\hsize=0.8\hsize\raggedright\arraybackslash}X|>{\hsize=1.1\hsize\raggedright\arraybackslash}X|>{\raggedright\arraybackslash}p{2.2cm}|>{\raggedright\arraybackslash}p{1.35cm}|}
    \caption[Schutzmechanismen mit Auslegungswerten]{Schutzmechanismen mit Gefährdung, Auslegungswert und Geltungsbereich, nach Wirkprinzip geordnet. Es sind Softwarefunktionen ohne Sicherheitsbewertung.}
    \label{tab:schutzmechanismen}\\\hline
    \tabkopfzeile \tabkopf{Mechanismus} & \tabkopf{Gefährdung} & \tabkopf{Auslegungswert} & \tabkopf{Sim/Real} & \tabkopf{Abschn.}\\\hline
    \endfirsthead
    \multicolumn{5}{l}{\small\emph{Fortsetzung von Tabelle~\ref{tab:schutzmechanismen}}}\\\hline
    \tabkopfzeile \tabkopf{Mechanismus} & \tabkopf{Gefährdung} & \tabkopf{Auslegungswert} & \tabkopf{Sim/Real} & \tabkopf{Abschn.}\\\hline
    \endhead
    \multicolumn{5}{|l|}{\cellcolor{baHellgrau}\textbf{Begrenzen}}\\\hline
    Armtempo: Gelenk, Handpunkt und Ellbogen & Stoß, Quetschen & 1,5\,rad/s; 0,25\,m/s & real; Sim 1,5\,rad/s, 0,5\,m/s & \ref{sec:regelkette}\\\hline
    Anfahrrampe, Schleppfehlerbegrenzung & Ruck, davonlaufender Sollwert & 3,0\,rad/s$^2$; 0,3\,rad & nur real & \ref{sec:regelkette}\\\hline
    Obergrenzen der Reglerverstärkung & zu harte Gelenke & Arm kp 80, kd 3; Handgelenk kp 50, kd 3; Taille kp 250, kd 6 & nur real & \ref{sec:regelkette}\\\hline
    Gewichtsrampe bei Übernahme und Rückgabe & Sprung beim Umschalten & 2,0\,s, mindestens 1\,s & real; Sim 0\,s & \ref{sec:regelkette}\\\hline
    Gehgeschwindigkeit & zu schnelles Loslaufen & vorwärts 0,4\,m/s, seitlich 0,3\,m/s, Drehen 0,4\,rad/s & real; Sim: Bereich der Policy & \ref{sec:onboard_anbindung}\\\hline
    \multicolumn{5}{|l|}{\cellcolor{baHellgrau}\textbf{Vermeiden}}\\\hline
    Kollisionsprüfung zur Laufzeit (Umgebung, Selbstkollision) & Kollision mit Umgebung oder Körper & Untergrenze Hand 0\,cm, Arm 1\,cm; langsamer unter 5\,cm; Selbstkollision Marge 3\,cm & beide & \ref{sec:kollisionsmodell}\\\hline
    Geplante Bahnen für gespeicherte Posen & Kollision auf dem Weg & Abstand 5\,cm in der Bahnmitte, mindestens 1,5\,cm; Trichter 15\,cm & beide & \ref{sec:kollisionsmodell}\\\hline
    Freigabe des Gehens: Bodenkontakt, Verriegelung & Loslaufen ohne Bodenkontakt oder mit verdrehter Taille & Knielast mindestens 10\,Nm & Bodenkontakt real; Verriegelung beide & \ref{sec:onboard_anbindung}\\\hline
    Zweite Bestätigung (Startleiste), Schieber am Smartphone & unbeabsichtigtes Dämpfen oder Loslaufen & zweites Tippen innerhalb 4\,s & Startleiste real; Schieber beide & \ref{sec:bedienwege}\\\hline
    Sicherheitsbestätigung vor dem Start, kein selbsttätiger Start & unbeabsichtigte Bewegung & -- & nur real & \ref{sec:bereitstellung}, \ref{sec:bedienwege}\\\hline
    Geplante Rückgabe der Arme, auch bei Programmende & Sprung bei der Rückgabe & Zeitgrenze 12\,s; Programmende höchstens 2,5\,s mit 0,8\,rad/s & nur real & \ref{sec:regelkette}\\\hline
    \multicolumn{5}{|l|}{\cellcolor{baHellgrau}\textbf{Überwachen}}\\\hline
    Totmann am Gamepad für das Gehen & Weiterlaufen bei Ausfall der Bedienperson & Taste mindestens halb gedrückt & beide & \ref{sec:bedienwege}\\\hline
    Zeitüberwachung der Steuersignale & Ausfall von Verbindung oder Bedienperson & Gehen 0,3\,s, Hand 0,5\,s ohne Gamepad-Signal; Funkprüfung 0,3\,s; Gehbefehl 0,5\,s, gültig 1,0\,s & Gamepad beide; Gehbefehl real & \ref{sec:onboard_anbindung}, \ref{sec:bedienwege}\\\hline
    Statusalter und Herzschlag am Smartphone & veraltete Daten, Funkausfall & Sperre nach 1,5\,s ohne Status; Steuerhoheit 2,0\,s, Sitzung 4,0\,s & beide & \ref{sec:bedienwege}\\\hline
    Sturzerkennung & Weiterbetrieb nach Sturz & Neigung über etwa 60$^\circ$ länger als 0,3\,s & nur Sim & \ref{sec:ersatzregler}\\\hline
    \multicolumn{5}{|l|}{\cellcolor{baHellgrau}\textbf{Anhalten}}\\\hline
    Halt-Stufen (Tabelle~\ref{tab:haltstufen}), Not-Halt auf jedem Bedienweg, Quittieren nicht am Gamepad & Weiterfahren bei Gefahr, unbeabsichtigtes Wiederanlaufen & Voll-Notaus: 2,0\,s halten & beide; Beine verschieden & \ref{sec:onboard_anbindung}, \ref{sec:bedienwege}\\\hline
\end{xltabular}
\endgroup
% Fundstellen je Zeile (Doku UND Code; bei Abweichung gilt der Code, Abweichungen stehen im Bericht):
% 1 Arm-Tempo:
%     Doku: docs/04-T-Umgebungsvariablen.md "Arme (echter Roboter)": G1_ARM_VEL_LIMIT 1.5 rad/s, G1_EE_VEL_LIMIT 0.25 m/s; docs/50-A-Echter-Roboter.md E "Harte Grenzen nur auf real ... Gelenkgeschwindigkeit <= 1,5 rad/s, Hand <= 0,25 m/s";
%           docs/20-T-Arm-Manipulation.md main_loop Schritt 6-7 ("_limit_cartesian_speed, überwacht Hand-TCPs + Ellbogen").
%     Code: manipulation/arm_controller.py:165 (arm_velocity_limit 1.5), :171 (ee_velocity_limit 0.25), :1123-1124 (REAL_MAX_ARM_VEL 1.5, REAL_MAX_EE_VEL 0.25), :1140-1147 (Korrektur auf Obergrenze),
%           :555-561 (überwachte Punkte: Hand-TCPs und Ellbogen), :1584 (_limit_cartesian_speed); launch/bringup_real.launch.py:106-107; Simulation: launch/bringup_sim.launch.py:112 (ee_velocity_limit '0.5'),
%           arm_velocity_limit in der Sim = Node-Default 1.5 (arm_controller.py:165, nicht durchgesetzt, da _apply_real_safety_limits nur real :1029).
% 2 Anfahrrampe/Schleppfehler:
%     Doku: docs/04-T: G1_ARM_ACCEL_LIMIT 3.0 rad/s², G1_ARM_MAX_TRACKING_ERROR 0.3 rad; docs/50-A E "Sanftes Anfahren und begrenzter Schleppfehler (real)"; docs/20-T Konfiguration (arm_accel_limit, max_tracking_error "0 = aus").
%     Code: arm_controller.py:330 (max_tracking_error 0.0 = aus), :336 (arm_accel_limit 0.0 = aus), :1403-1418 (_limit_accel, begrenzt nur Schneller-Werden), :1419-1430 (_limit_tracking_error, nur use_robot und Gewicht "on");
%           launch/bringup_real.launch.py:110 ('0.3'), :111 ('3.0'); Sim: manipulation_launcher.launch.py:78-79 (Default 0.0), bringup_sim setzt nichts -> aus.
% 3 Reglerverstärkung:
%     Doku: docs/50-A E "Harte Grenzen nur auf real: Gains höchstens kp 80/kd 3 (Handgelenk kp 50)"; "Obergrenze 250 / 6" für die Taille (docs/50-A E "Taille wird mitgehalten").
%     Code: arm_controller.py:1119-1120 (REAL_MAX_GAINS kp_low 80, kd_low 3, kp_wrist 50, kd_wrist 3, kp_waist 250, kd_waist 6), :1131-1136 (Korrektur mit Fehlermeldung).
%     Abweichung (Ergänzung, kein Widerspruch): Handgelenk-kd-Obergrenze 3 steht nur im Code. Eingestellt sind real kp 60/kd 1,5 (Arm) und kp 40/kd 1,5 (Handgelenk), launch/bringup_real.launch.py:115-118;
%           Sim 150/12 und 40/4: manipulation_launcher.launch.py:53-56. Tabelle nennt die Obergrenzen als Auslegungswert.
% 4 Gewichtsrampe:
%     Doku: docs/20-T "rt/arm_sdk-Gewichtsrampe" (Sim 0.0, real 2.0 s), docs/50-A E "rampt das arm_sdk-Gewicht 0→1 über 2,0 s", "Gewichtsrampen ≥ 1 s".
%     Code: arm_controller.py:310-311 (arm_weight_ramp_up_s/down_s 2.0), :1122 (REAL_MIN_RAMP_S 1.0), :1137-1139 (Korrektur), launch/bringup_sim.launch.py:109-110 ('0.0').
% 5 Gehgeschwindigkeit:
%     Doku: docs/04-T "Gehen (echter Roboter)": G1_MAX_VX/VY/VYAW 0.4/0.3/0.4; docs/50-A C "Geh-Limits Default vorwärts 0,4 m/s, seitlich 0,3 m/s, drehen 0,4 rad/s".
%     Code: navigation/loco_client.py:145-147 (max_vx 0.4, max_vy 0.3, max_vyaw 0.4), :731-733 (Skalierung), launch/bringup_real.launch.py:60-62, :174-176 (durchgereicht).
%           Simulation: loco_sim.py:559-562 (walk_policy.scale_command, keine eigene Grenze); Bereiche der Lauf-Policies policies/g1_wholebody/deploy.yaml:12 (lin_vel_x [-0.5, 1.0]),
%           agile_velocity_g1/deploy.yaml:66 ([-0.5, 0.5]), agile_velocity_height_g1/deploy.yaml:104 ([-0.5, 1.0]) -> "bis 1,0 m/s vorwärts".
% 6 Kollisionsprüfung zur Laufzeit:
%     Doku: docs/20-T "Laufzeit-Gate (_apply_collision_gate)" (Hand gate_hand_floor 0 = aufsetzen erlaubt, Arm gate_arm_floor 1 cm, gate_slow_dist 5 cm, Selbstkollision Margenband 3 cm), Konfigurationstabelle.
%     Code: arm_controller.py:202-206 (gate_hand_floor 0.0, gate_arm_floor 0.01, gate_slow_dist 0.05), :1624-1632 (Untergrenzen), :1714-1716 (Verlangsamung bis 25 %); utils/ik_solver.py:53 (collision_margin=0.03), :89, :293.
%           Hand gegen Greif-Objekte keine Untergrenze (docs/20-T, arm_controller.py:1624-1625).
% 7 Geplante Bahnen:
%     Doku: docs/20-T "arm_planner.py" (Abstand in der Bahnmitte clearance 5 cm, Untergrenze 1,5 cm, relax_radius 15 cm), Konfigurationstabelle planner_*.
%     Code: manipulation/arm_planner.py:72-82 (DEFAULTS clearance 0.05, min_clearance 0.015, relax_radius 0.15), arm_controller.py:190-195 (Parameter).
% 8 Freigabe des Gehens:
%     Doku: docs/04-T G1_GROUND_CHECK (knee_load), G1_KNEE_LOAD_MIN_NM 10.0; docs/30-T "Bodenkontakt vor »Walk«"; docs/50-A E "WALK ist gesperrt, bis balanciert"; docs/24-T-Hueftdrehung.md (Gehen nur mit Hüfte in der Mitte).
%     Code: navigation/loco_client.py:190 (ground_check 'knee_load'), :192 (knee_load_min_nm 10.0), :334-341 (_ground_check), :726-728 (Move nur balanciert), :426-435 (Gehen mit 0, solange Hüfte nicht in der Mitte);
%           launch/bringup_real.launch.py:187-188; Simulation: navigation/loco_sim.py:454-456 (wartet auf walk_ok der Hüfte), kein Bodenkontakt-Test im Ersatzregler.
% 9 Zweite Bestätigung:
%     Doku: docs/10-A-Demo-GUI.md "Startleiste (nur echter G1): Damp und ein Walk ... brauchen einen zweiten Tipp innerhalb weniger Sekunden"; docs/12-A-Handy.md "Schieber ganz nach rechts".
%     Code: teleoperation/demo_gui.py:141 (QUIT_CONFIRM_MS = 4000), :1520-1531 (zweites Tippen, Timer); teleoperation/phone/web/app.js:490-515 (Bestätigung per Schieber).
% 10 Sicherheitsbestätigung/kein selbsttätiger Start:
%     Doku: docs/04-T G1_REAL_CONFIRM ("ohne startet REAL mit --yes nicht"); docs/50-A E "Kein Auto-Start" und C ("Haken Sicherheitshinweise gelesen"); docs/40-T "Auto-Start (GREIFEN)" nur Sim.
%     Code: start.sh:197-198, g1_gui.py:1143; teleoperation/sim_autostart.py (nur is_sim_mode).
% 11 Geplante Rückgabe:
%     Doku: docs/20-T "Übergabe an den Onboard-Regler" (release_timeout_s 12 s, shutdown_return_s 2,5 s, ≤ 0,8 rad/s), docs/50-A E "Zurückgeben ohne Schwung", "Programmende".
%     Code: arm_controller.py:319-324 (release_to_onboard_pose True, release_tolerance 0.05, release_timeout_s 12.0, shutdown_return_s 2.5), :1352-1354 (nur bei ramp > 0 -> real), :1505 (vmax = min(arm_velocity_limit, 0.8)), :3625 (Zeitgrenze).
% 12 Totmann:
%     Doku: docs/11-A-Controller.md ("Gelaufen wird nur, solange L2 gehalten wird (Totmann)"); docs/51-A-Notaus.md.
%     Code: teleoperation/ps4_walk.py:47 (DEADMAN_THRESHOLD = 0.5, "L2 mindestens halb gedrückt"), :52.
% 13 Zeitüberwachung der Steuersignale:
%     Doku: docs/51-T-Notaus.md "Zeitgrenzen ohne Notaus-Taste" (ps4_walk 0,3 s; loco_client 0,5 s; G1-Firmware Move() gilt 1 s), docs/11-T-Controller.md "Sicherheit" (joy_timeout_s 0.5, stall_timeout_s 0.3).
%     Code: teleoperation/ps4_walk.py:230 (joy_timeout_s 0.3), ps4_arm_teleop.py:751 (joy_timeout_s 0.5), ps4_joystick.py:68 (stall_timeout_s 0.3), navigation/loco_client.py:148 (cmd_vel_timeout 0.5), launch/bringup_real.launch.py:177,
%           loco_client.py:160 + :428 (move_duration_s 1.0, von der Umgebung als Befehlsdauer an den Roboter übergeben).
%     Abweichung (Zuordnung): Doku schreibt die Gültigkeit von 1 s der Firmware zu; im Code setzt die Umgebung die Dauer selbst (Parameter move_duration_s, Standard 1,0 s, nicht überschrieben). Wert stimmt, Tabelle nennt "Gültigkeit am Roboter 1,0 s".
%     Abweichung (Ergänzung): Doku 51-T nennt nur die 0,3 s für ps4_walk; die Hand-Steuerung hält nach 0,5 s an (ps4_arm_teleop.py:751). Beides in der Tabelle.
% 14 Smartphone:
%     Doku: docs/12-T-Handy.md "Zeiten" (Herzschlag 5 Hz, Steuerhoheit 2 s, Sitzung 4 s, Status 5 Hz, "Daten alt" nach 1,5 s), docs/12-A-Handy.md Ereignistabelle.
%     Code: teleoperation/phone/web/app.js:12 (FRESH_MS = 1500), phone/gateway.py:47-48 (CONTROL_TIMEOUT_S 2.0, SESSION_TIMEOUT_S 4.0); Sperre außer Not-Halt: gateway.py:60 ("estop" ohne Steuerhoheit/Statusbezug).
% 15 Sturzerkennung:
%     Doku: docs/30-T "Zustandsmaschine" ("Sturz-Erkennung (_fallen, IMU-Neigung über fall_gz für fall_debounce_s) → automatisch DAMP").
%     Code: navigation/loco_sim.py:144-147 (fall_gz -0.5 = "Neigung ~>60 Grad", fall_debounce_s 0.3), :566-575 (_fallen). Im Realbetrieb keine solche Funktion in loco_client.py.
% 16 Halt-Stufen:
%     Doku: docs/51-A-Notaus.md, docs/51-T-Notaus.md; Code: teleoperation/ps4_arm_teleop.py:736 (full_estop_hold_s 2.0); Not-Halt auf jedem Bedienweg: Demo-GUI, ps4_arm_teleop (PS), Smartphone (gateway.py:60), siehe docs/51-T-Notaus.md "Topics";
%           Quittieren am Gamepad nicht möglich: docs/51-A-Notaus.md Tabelle "Quittieren", ps4_arm_teleop.py (Sperre bis /g1pilot/start).
%     Abweichung (nur Doku): docs/20-T "Bekannte Einschränkungen" behauptet, die Taille werde vom arm_controller "bewusst nie beschrieben"; der Code hält die Taille real (waist_hold.py, arm_controller.py:1174-1180,
%           estop_waist_hold :227) und docs/50-A E, 51-T, 24-T beschreiben das so. Tabelle folgt dem Code.


% ----------------------------------------------------------------------
% Tabelle 5: Bedienwege (Abschnitt 6.4.2); ersetzt die Entwurfstabelle in Kapitel/07_Umsetzung.tex
% ----------------------------------------------------------------------
\begin{table}[htbp]
    \tabstil
    \setlength{\tabcolsep}{3pt}
    \centering
    \caption[Bedienwege und ihre Funktionen]{Bedienwege und ihre Funktionen zum Stand der Umgebung. \enquote{ja} bedeutet, dass der Bedienweg die Funktion selbst auslöst; die Logik liegt in allen Fällen in den Funktionsknoten.}
    \label{tab:bedienwege}
    \begin{tabularx}{\textwidth}{|>{\raggedright\arraybackslash}p{2.0cm}|L|L|L|L|L|}\hline
        \tabkopfzeile & \tabkopf{Bedienoberfläche} & \tabkopf{Gamepad} & \tabkopf{Smartphone} & \tabkopf{Zielmarker in RViz} & \tabkopf{Arm-API}\\\hline
        Einsatz \pruefen{Zielgruppe je Bedienweg von Nils bestätigen} & Standard, Vorführung & am Roboter & Zusatz zur Bedienoberfläche & Entwicklung & Programme ohne ROS\\\hline
        Arme & ja (Hand verfahren, Grundstellung) & ja (Hand verfahren, Arm spiegeln) & nur über Posen & ja (Zielmarker) & ja (Gelenkwinkel, Zielpose)\\\hline
        Hände & ja & ja & ja & nein & ja\\\hline
        Gehen & ja (Totmann); Navigation nur Sim & ja (Totmann) & nein & Navigationsziel (nur Sim) & nein\\\hline
        Posen und Abläufe & ja (anfahren, speichern, Abläufe) & Ablauf weiter, Stopp & ja; Abläufe zusammenstellen nur am PC & nein & nur Speichern\\\hline
        Not-Halt & ja & ja & ja, auch in der Rolle Ansicht & nein & nein (nur Abbrechen)\\\hline
        Quittieren & ja & nein (absichtlich) & ja (Rolle Bediener) & nein & nein\\\hline
        Simulation / real & beide & beide & beide & beide & beide\\\hline
    \end{tabularx}
\end{table}
% Fundstellen je Zeile (Doku / Code); Bedienwege nur soweit zum Stand 25d5c1a vorhanden. Entfernt seit 30.10.2026: Stream-Deck-Oberfläche (NOTICE.md); alter Joystick-Modus (docs/13) ist ein Rückfall,
% kein Bedienweg des Konzepts (Plan 6.4.2 nennt ihn nicht).
% Zielgruppe:   Doku docs/1-Bedienung.md Tabelle "Möglichkeit/Wofür" (Demo-GUI "Standard am PC und für Vorführungen"; Controller "direkt am Roboter"; Handy "Zusatz zur Demo-GUI");
%               docs/10-A-Demo-GUI.md (Messe, Besuch, Lehre); docs/20-A-Arm-Manipulation.md (RViz); docs/21-A-Arm-API.md ("eigenes, unabhängiges Projekt ... ohne ROS 2").
%               Eine Zielgruppe (Laie/Entwickler) ist für Gamepad und Smartphone nirgends angegeben -> \pruefen. "Entwicklung" für RViz folgt Plan 6.4.2 ("Zielmarker in RViz (Entwickler)").
% Arme:         Doku docs/10-A-Demo-GUI.md (GREIFEN: "Hand bewegen", Grundstellung, Sichere Pose), docs/11-A-Controller.md Belegung (Sticks, Share 1 s = Grundstellung, Arm spiegeln),
%               docs/12-A-Handy.md ("Gehen und Arm-/Hand-Jog ... Das Handy kann das bewusst nicht"), docs/20-A-Arm-Manipulation.md (Marker), docs/21-A-Arm-API.md (/arm/joints, /arm/pose).
%               Code manipulation/interactive_marker.py (Marker), teleoperation/phone/gateway.py:58-75 (ACTIONS: home, safe, pose, move; kein Jog), manipulation/arm_api.py:292.
% Hände:        Doku docs/10-T-Demo-GUI.md (Seite HÄNDE: Finger, Kraft, Kraftzonen), docs/11-A (Kreis/Viereck), docs/12-T "hand" (öffnen/schließen), docs/21-A ("Optional: hands (6 Fingerwinkel je Hand ...)"); RViz-Marker senden nur Handziele.
%               Code manipulation/arm_command.py:80, :175, :227-231 (hands im Wire-Format), gateway.py:67 ("hand").
% Gehen:        Doku docs/10-A-Demo-GUI.md "GEHEN" (Knopf, Pfeiltasten, Tot-Mann-Prinzip, AUTO NAV), docs/11-A "Gehen mit L2", docs/12-A "Was das Handy nicht kann: Gehen",
%               docs/31-A-Navigation.md ("In RViz: 2D Goal Pose", "Die Navigation läuft nur in der Simulation").
%               Abweichung zum Entwurf in Kap. 7: dort "Rviz: Beine = Auto-Nav" ohne Hinweis auf "nur Simulation"; Arm-API ohne Beine (stimmt); Smartphone: Modus-/Startleistenwahl (gateway.py:65-66 "loco", "mode") ist kein Gehen.
% Posen/Abläufe: Doku docs/12-A-Handy.md "Posen" (anfahren, speichern, Abläufe; "Zusammenstellen geht nur am PC"), docs/11-A (L3 = Ablauf weiter; Kreuz = Stopp), docs/21-A (/arm/save; kein Anfahren gespeicherter Posen über HTTP),
%               docs/23-T-Posen-und-Ablaeufe.md. Code arm_api.py:289-291 (nur /arm/save), gateway.py:68-74.
%               Abweichung zum Entwurf in Kap. 7: dort "Arm-API: Posenspeicher x" -- die API speichert nur, sie fährt keine Posen an.
% Not-Halt:     Doku docs/51-T-Notaus.md "Topics" (Sender: Demo-GUI, ps4_arm_teleop, Handy, alter Joystick); docs/12-T "estop" (Rolle Ansicht); docs/21-A (/arm/cancel nur Abbrechen).
%               Code gateway.py:60, link.py:68; arm_api.py:285 (/arm/cancel).
%               Abweichung zum Entwurf in Kap. 7: Zeile Not-Halt fehlte dort.
% Quittieren:   Doku docs/51-A-Notaus.md Tabelle "Quittieren" (Demo-GUI, Handy mit Recht »Bedienen«, Controller "geht nicht, absichtlich"); Code gateway.py:64 ("ack": Bediener, Steuerhoheit).
% Sim/real:     Doku docs/1-Bedienung.md ("Alle Wege senden dieselben Befehle"), docs/20-A ("RViz: Sim optional, Real standardmäßig an"), docs/12-A ("G1 REAL" / "G1 SIM"), docs/31-A (Navigation nur Sim).


% ----------------------------------------------------------------------
% Tabelle 6: Andockpunkte für externe Projekte (Abschnitt 6.5.1)
% Rein konzeptionell; bisher ist kein externes Projekt angebunden (Nils 16.11., N10).
% ----------------------------------------------------------------------
\begin{table}[htbp]
    \tabstil
    \setlength{\tabcolsep}{3pt}
    \centering
    \caption[Andockpunkte für externe Projekte]{Andockpunkte für externe Projekte. Je tiefer der Andockpunkt liegt, desto größer ist die Freiheit und desto weniger Schutz wirkt.}
    \label{tab:andockpunkte}
    \begin{tabularx}{\textwidth}{|>{\raggedright\arraybackslash}p{2.6cm}|L|L|>{\raggedright\arraybackslash}p{1.7cm}|>{\raggedright\arraybackslash}p{2.6cm}|}\hline
        \tabkopfzeile \tabkopf{Andockpunkt} & \tabkopf{Projekt liefert} & \tabkopf{Schutz wirkt} & \tabkopf{Sim / real} & \tabkopf{Detail in}\\\hline
        Daten (Szenen, Hand-Profile, Lauf-Policies, Posen und Abläufe) & Dateien nach festem Format, keine Programmierung & ja, Daten gehen in Planung und Prüfung ein & Sim; Posen beide & Abschnitte~\ref{sec:szenen}, \ref{sec:haende}, \ref{sec:ersatzregler}, Anhang~\ref{anh:dokumentation}\\\hline
        HTTP-Schnittstelle der Arme & Zielposen und Gelenkwinkel als JSON & ja, Prüfungen im Armregler & beide; nur lokal & Abschnitt~\ref{sec:entwickler_api}\\\hline
        ROS-2-Kanäle & Ziele und Zustände auf den Kanälen der Bedienwege & Schutz der Funktionsknoten wirkt; Prüfungen einzelner Bedienwege (etwa Rollen und Steuerhoheit des Smartphones) nicht & beide & Abschnitt~\ref{sec:navigation}, Anhang~\ref{anh:dokumentation}\\\hline
        Herstellerschnittstelle direkt & eigenes Programm über das Unitree-SDK & nein, kein Schutz der Umgebung & beide & Abschnitt~\ref{sec:bridge}, Anhang~\ref{anh:dokumentation}\\\hline
    \end{tabularx}
\end{table}
% Fundstellen je Zeile (Doku / Code):
% Szenen:       Doku unitree_mujoco/scene_editor/README.md ("Hindernis vs. Greif-Objekt ... Namenskonvention Präfix grasp_"; Editor "Objekte per Maus platzieren"), docs/20-T "Kugelmodell" (Objekte aus /scene_markers als Boxen, Planer + Laufzeit-Gate),
%               docs/40-T "Umgebungen erzeugen". Code navigation/scene_bridge.py:190 (obstacles + grasp), manipulation/arm_geometry.py (Environment, env_dist).
%               "nur Simulation": scene_bridge ist laut docs/02-T "(Sim)"; Szenen sind MuJoCo-Modelle.
% Hand-Profile: Doku docs/25-T-Eigene-Haende.md ("Ordner hands/<name>/ mit hand.yaml und einer STL ... Eigene Profile wirken nur bei G1_SIM_MODE=true"; "Kollision ... Gate und Planer-Kugeln aus denselben Teilen"); docs/04-T G1_HANDS.
%               Code utils/hand_profile.py, utils/ik_solver.py (_GATE_ARM_KEYS enthält hand_tool, laut docs/25-T).
% Lauf-Policies: Doku docs/30-T "Beteiligte Dateien" (policies/<name>/ mit deploy.yaml, ONNX), "Abfangen", "Zustandsmaschine" (Sturz -> DAMP); Code navigation/walk_policy.py (load_walk_policy am Feld format der deploy.yaml),
%               loco_sim.py:144-147, :566-575.  "Arm-Schutz ja": die Policy betrifft nur die Beine, die Arme laufen weiter über den Armregler.
% Posen/Abläufe: Doku docs/23-T ("Jeder Schritt ... als /g1pilot/arm_command type=pose -- geplant, mit Kollisionsschutz"), docs/23-A; Code teleoperation/sequence_runner.py, manipulation/pose_store.py (Format Version 3).
% HTTP:         Doku docs/21-T-Arm-API.md "Sicherheit" (Default-Bind 127.0.0.1, "Die Brücke kann sie prinzipbedingt nicht umgehen"), docs/21-A "Endpunkte"; Code manipulation/arm_api.py:82 (bind_host "127.0.0.1"), :85 (auth_token leer),
%               launch/manipulation_launcher.launch.py:96-100 (arm_api_host 127.0.0.1, arm_api_port 8770, arm_api_token ""), arm_command.py (Format-Prüfung, Gates im arm_controller).
% ROS-2-Kanäle: Doku docs/31-A-Navigation.md ("virtueller Joystick — denselben Signalweg, den auch eine manuelle Bedienung nutzt"), docs/31-T "Warum der Umweg über einen Fake-Joystick?", docs/05-T-Nodes-und-Topics.md.
%               Code navigation/joy_to_cmdvel.py:4-12.
% Hersteller:   Doku docs/21-T-Arm-API.md "Ohne eigenen HTTP-Client: direkt über DDS" ("Es genügt, die eigene DDS-Verbindung auf Domain 1, Interface lo zu zeigen ... umgeht den kompletten g1pilot-ROS-Graph inklusive Arm-Controller und Sicherheitsgates");
%               Code utils/common.py:21-23, :36-40. Entwurf im Plan sagte "ohne Änderung"; die Doku verlangt das Umstellen von Domain und Schnittstelle, Tabelle folgt der Doku.
% Hinweis: Es ist bisher kein externes Projekt angebunden (plan-kap6.md N10); die Tabelle behauptet keine erfolgte Anbindung.

% Hinweis (Koordinator, Textauftrag 6.5): genau vier Zeilen; Smartphone-Rolle heißt im Code "Bediener". Die Zeile ROS-2-Kanäle stützt sich auf teleoperation/phone (Rollen/Steuerhoheit) und die Funktionsknoten (Gates im arm_controller, loco_client); Zelle bei Bedarf von Nils prüfen lassen.
